% Depth-partition table: content depth vs. zero-shot readable depth vs. gap.
% content-j = truncation-downstream probe knee98 (~0.45L, near scale-invariant);
% readout-crash = zero-shot QCMem single-recall 50% collapse point on RULER
% niah_single 16k (bracket sweep); gap = content - readout = adapter's job.
% Numbers from status/QCMEM_J_DETERMINATION.md (probe + readout bracket).
\begin{table}[t]
\centering
\small
\setlength{\tabcolsep}{5pt}
\begin{tabular}{@{}lcccc@{}}
\toprule
\textbf{Model} & \textbf{$L$} & \textbf{content-$j$} & \textbf{readout-crash} & \textbf{gap} \\
 & & ($/L$) & ($/L$) & ($\Delta/L$) \\
\midrule
0.6B & 28 & 13 ($0.48$) & $\sim$2.6 ($0.09$) & \textbf{0.39} \\
1.7B & 28 & 13 ($0.48$) & $\sim$6 ($0.22$)  & 0.26 \\
4B   & 36 & 16 ($0.44$) & $\sim$11 ($0.31$) & 0.13 \\
8B   & 36 & 16 ($0.44$) & $\sim$11 ($0.30$) & 0.14 \\
14B  & 40 & 18 ($0.46$) & $\sim$15 ($0.375$) & 0.085 \\
32B  & 64 & 27 ($0.42$) & $>$27 ($>0.42$)   & \textbf{$\approx$0} \\
\bottomrule
\end{tabular}
\caption{\textbf{Depth partition: content vs.\ readable depth.} The
\emph{content-$j$} (semantic peak, from a truncation-downstream linear probe)
sits at $\approx$$0.45L$ and is \emph{near scale-invariant} across two orders of
magnitude. The zero-shot \emph{readout-crash} depth (where CoMem's zero-shot
single recall drops below 50\% on RULER \texttt{niah\_single} 16k) instead
\emph{deepens monotonically with scale}, from $0.09L$ (0.6B) to $>0.42L$ (32B,
no crash observed). Their difference is the \textbf{gap} the self-distilled
adapter must close: it is largest for tiny models (0.39$L$ at 0.6B) and vanishes
by 32B, where the zero-shot readout already reaches the content peak. Hence
small models benefit most from the adapter, while 32B needs it least.}
\label{tab:depth}
\end{table}
